\documentclass[11pt]{article}

\usepackage[margin=25mm]{geometry}

\usepackage{amsmath,amssymb,fontspec,unicode-math}

\IfFontExistsTF{DejaVu Serif}{\setmainfont{DejaVu Serif}}{\setmainfont{Latin Modern Roman}}

\IfFontExistsTF{Latin Modern Math}{\setmathfont{Latin Modern Math}}{\setmathfont{DejaVu Math TeX Gyre}}

\usepackage{xurl,hyperref}

\setlength{\parindent}{0pt}\setlength{\parskip}{6pt}

\begin{document}

\title{Feature Distillation is the Better Choice for Model-Heterogeneous Federated Learning}

\author{ELVIA preliminary AI analysis}\date{21 October 2026}\maketitle

Preliminary assessment: Minor.\par

Independent scientific review remains pending. Author adoption has not been requested.

Rectangular orthonormal projection preserves only retained-subspace geometry, and separate projection orthogonality does not prevent conflicting updates to shared features. The square exponential identity, dimensional alignment and empirical regularization benefits remain compatible with these limits.

\section*{Review scope}

Claim-focused reading of the original paper and relevant author-released source, with direct mathematical and static-source checks. No author program, historical verifier, model training or published benchmark was executed.

\section*{Problem 1: Dimension reduction cannot preserve all source geometry}

Source: Section 3.3, Eqs.(6)–(8), pp.6–7

Assessment: supported; impact: minor

\subsection*{Source claim (paraphrase)}

The text invokes rotational invariance after truncating a square orthogonal matrix to align a lower-dimensional teacher, attributing feature-shape preservation to that map.

\subsection*{Analysis}

In a favorable orientation let \(A=Q_d^{\top}\), \(AA^{\top}=I_d\), \(P=A^{\top}A\). Then \(\|Av\|^2=\|v\|^2-\|(I-P)v\|^2\). With \(D=4,d=2,Q=\exp(0)=I_4\), the projection retains only the first two coordinates: \(e_3,e_4\) collapse to zero despite distance \(\sqrt2\). The data matrix containing the standard basis has full row rank, so a full-rank assumption does not rescue preservation. Multiplication by any fixed projection remains linear regardless of sample count/rank.

\subsection*{Finding and scope}

Norms/distances survive exactly on retained directions, not all source features. Tensor-dimension alignment and possible task-relevant approximate preservation remain intact.

\subsection*{Proposed correction}

State contraction and retained-subspace isometry. If discarded pairwise components are at most \(\varepsilon\) times original distance, retained distance lies between \(\sqrt{1-\varepsilon^2}\) and one times original distance; that data-dependent condition is not established here.

A proposed correction is described; author adoption is not established.

\section*{Problem 2: Separate orthogonality does not eliminate shared-gradient conflicts}

Source: Section 3.3, conflict-prevention explanation and Eq.(9)

Assessment: supported; impact: minor

\subsection*{Source claim (paraphrase)}

Per-map identities \(A_iA_i^{\top}=I\) do not control cross-map inner products \(g_i^{\top}g_j=r_i^{\top}A_iA_j^{\top}r_j\).

\subsection*{Analysis}

Use coordinate projections from \(\mathbb R^4\) retaining first \(2\) and first \(3\) coordinates, both permitted truncations of \(I_4\). At \(z=(1,1,1,0)\), apply softmax to projected features and use teachers \((0.9,0.1)\), \((0.1,0.8,0.1)\). For teacher-to-student KL, shared gradients are \((-2/5,2/5,0,0)\) and \((7/30,-7/15,7/30,0)\), inner product \(-7/25<0\). For the printed reversed KL direction it is \(-\log9\log8/12<0\). Thus improving one term can worsen the other to first order under either convention. Softmax is an explicit favorable domain assumption, not recovered unavailable code.

\subsection*{Finding and scope}

Individual projection orthogonality permits this operational conflict. Conditioning improvements, learned conflict reduction and empirical benefit remain possible; no universal parameter-gradient result follows.

\subsection*{Proposed correction}

Describe a per-projection constraint and empirical regularizer; measure defined gradient conflicts. Cross-map orthogonality would suffice at the shared-feature level but is not imposed and may be infeasible; it is not offered as a validated replacement.

A proposed correction is described; author adoption is not established.

\section*{Limits}

Preliminary AI assessment; independent human scientific adjudication remains pending.\par

Analytic examples and release discrepancies do not remeasure the paper’s reported experiments.\par

\section*{Source identity}

Current official NeurIPS 2025 publisher PDF, freshly retrieved 21 October 2026

\url{https://proceedings.nips.cc/paper_files/paper/2025/file/96dd63ab8fd79163a5399e0b385f208d-Paper-Conference.pdf}

Reviewed PDF SHA-256: \nolinkurl{e9d916896e62a83fb01f7979f79dc1d95c9cd4327d063ab0d175a7551623035a}

Fresh HTTP 200 download from the official publisher; actual retrieval timestamp and SHA-256 recorded. The content is byte-identical to the retained proceedings PDF. No file modification time is used as provenance.

\end{document}
